\subsection{\texorpdfstring{\(\mathcal{L}(\mathbf{E};\mathbf{F})\) 的有界子集（拟完备情形）}{\textbackslash mathcal\{L\}(\textbackslash mathbf\{E\};\textbackslash mathbf\{F\}) 的有界子集（拟完备情形）}}

定理 2. --- 设 E 是一个局部凸 Hausdorff 空间，F 是一个局部凸空间，S 是 E 的闭的、凸的、均衡的、有界的且半完备的子集的一个族（III, 第 7 页）。\(\mathcal{L}(\mathrm{E};\mathrm{F})\) 的每个简单有界子集 H 对 S-拓扑都是有界的。

设 \(A \in S\)。空间 \(E_A\) 于是是一个 Banach 空间（III, 第 8 页, 推论），因而是桶型的。另一方面，H 在 \(\mathcal{L}(E_A; F)\) 中的典范像是简单有界的，因而是等度连续的（III, 第 25 页, 定理 1）。于是，所有 \(u(x)\)（对 \(u \in H\) 和 \(x \in A\)）所成的集在 F 中有界，这就证明 H 对 \(S\) -拓扑有界。

推论 1. --- 设 E 是一个局部凸 Hausdorff 空间，F 是一个局部凸空间，S 是 E 的有界子集的一个族。若 E 是半完备的，则 \(\mathcal{L}(\mathrm{E};\mathrm{F})\) 的每个简单有界子集对 S-拓扑都是有界的。

只需在把 S 中的集换成它们的闭凸均衡包之后应用定理 2，因为这并不改变 S-拓扑。

当 E 是半完备的（例如拟完备的）时，我们可以谈论 \(\mathcal{L}(\mathrm{E};\mathrm{F})\) 的有界子集而不必指明 S-拓扑，因为只要 S 是 E 的覆盖，这些子集对所有 S-拓扑都是相同的。

推论 2. --- 每个半完备的有界型空间都是桶型的。

这种空间的对偶的每个简单有界子集都是强有界的（推论 1），因而是等度连续的（III, 第 22 页, 命题 10）。

推论 3. --- 设 E 是一个局部凸空间。E 的每个对 \(\sigma(E, E')\) 有界的子集都是有界的。

设 A 是 E 的一个子集。说 A 对 \(\sigma(\mathrm{E}, \mathrm{E}')\) 有界，意思是 E 上每个连续线性型都在 A 上有界；说 A 有界，意思是 E 上每个连续半范数都在 A 上有界。设 N 是 0 在 E 中的闭包，\(\pi\) 是从 E 到 E/N 上的典范映射。E 上的连续线性型是形如 \(f \circ \pi\) 的映射，其中 \(f \in (\mathrm{E}/\mathrm{N})'\)，且我们对 E 上的连续半范数有类似的刻画。把 E 换成 E/N、A 换成 \(\pi(\mathrm{A})\)，我们就可限于 E 为 Hausdorff 的情形。

设 \(\mathfrak{S}\) 是 \(E'\) 的等度连续子集的集；当给 \(E'\) 赋予拓扑 \(\sigma(E', E)\) 时，\(E\) 可与 \((E')_{\mathfrak{S}}'\) 等同（III, 第 19 页, 推论 1）。\(E'\) 的每个闭的等度连续子集对 \(\sigma(E', E)\) 是紧的（III, 第 17 页, 推论 2），从而对 \(\sigma(E', E)\) 是完备的。现在只需应用定理 2。

